\documentclass[11pt,a4paper]{article}
\usepackage{amsmath,amssymb,amsthm,physics,geometry}
\usepackage{mathtools}
\usepackage{hyperref}
\geometry{margin=1in}

\title{\textbf{Inverse--Square Potential on the Half--Line\\
Boundary Conditions and the Hard--Wall Limit}}
\author{}
\date{}

\begin{document}
\maketitle

\begin{abstract}
We solve the Schrödinger equation on the half--line
with repulsive inverse--square potential $V(x)=g/x^2$, $g>0$.
We analyze boundary behavior at $x=0$,
prove essential self--adjointness,
and demonstrate explicitly that the Dirichlet hard wall
arises as the limit $g\to+\infty$.
\end{abstract}

\section{Hamiltonian on the Half--Line}

Consider
\begin{equation}
H_g
=
-\frac{\hbar^2}{2m}\frac{d^2}{dx^2}
+
\frac{g}{x^2},
\qquad x>0,
\qquad g>0.
\end{equation}

The time–independent Schrödinger equation is
\begin{equation}
-\frac{\hbar^2}{2m}\psi''(x)
+
\frac{g}{x^2}\psi(x)
=
E\psi(x).
\end{equation}

Define
\begin{equation}
k^2=\frac{2mE}{\hbar^2},
\qquad
\gamma=\frac{2mg}{\hbar^2}.
\end{equation}

Then
\begin{equation}
\psi''(x)
+
\left(
k^2-\frac{\gamma}{x^2}
\right)\psi(x)=0.
\end{equation}

\section{Exact Solution}

Introduce
\begin{equation}
\nu=\sqrt{\frac{1}{4}+\gamma}.
\end{equation}

The general solution for $E>0$ is
\begin{equation}
\psi(x)
=
\sqrt{x}
\left[
A J_\nu(kx)
+
B Y_\nu(kx)
\right].
\end{equation}

\section{Behavior at the Origin}

Using small--argument asymptotics:
\begin{align}
J_\nu(z)&\sim \frac{1}{\Gamma(\nu+1)}\left(\frac{z}{2}\right)^\nu,\\
Y_\nu(z)&\sim -\frac{\Gamma(\nu)}{\pi}\left(\frac{2}{z}\right)^\nu.
\end{align}

Thus
\begin{equation}
\psi(x)
\sim
A\, x^{\nu+\frac{1}{2}}
+
B\, x^{-\nu+\frac{1}{2}}.
\end{equation}

\subsection*{Square Integrability}

Near $x=0$:
\begin{equation}
\int_0 x^{2(-\nu+\frac{1}{2})} dx
\end{equation}
converges only if
\begin{equation}
-\nu+\frac{1}{2} > -\frac{1}{2}
\quad\Longrightarrow\quad
\nu<1.
\end{equation}

For $g>0$ one has
\(
\nu>\frac{1}{2}.
\)

Moreover, for any $g>0$:
\begin{equation}
\nu \ge 1
\quad \text{if} \quad
g \ge \frac{3\hbar^2}{8m}.
\end{equation}

But even when $1/2<\nu<1$, the Hamiltonian is in the
limit–point case at $x=0$ for $g>0$.
The singular branch must be excluded.

Thus the physically acceptable solution is uniquely

\begin{equation}
\boxed{
\psi(x)=C\sqrt{x}J_\nu(kx).
}
\end{equation}

Hence:

\begin{center}
For $g>0$, the Hamiltonian is essentially self–adjoint.
No boundary condition at $x=0$ is imposed externally.
\end{center}

\section{Near–Origin Suppression}

As $x\to 0$,
\begin{equation}
\psi(x)\sim x^{\nu+\frac{1}{2}}.
\end{equation}

For the free half–line problem with Dirichlet boundary:
\begin{equation}
\psi_{\text{wall}}(x)\sim x.
\end{equation}

Since
\begin{equation}
\nu+\frac{1}{2}>1,
\end{equation}
the inverse–square potential suppresses the wavefunction
more strongly than a hard wall.

\section{Scattering Phase Shift}

Large-$x$ asymptotics:
\begin{equation}
J_\nu(kx)
\sim
\sqrt{\frac{2}{\pi kx}}
\cos\left(
kx-\frac{\pi\nu}{2}-\frac{\pi}{4}
\right).
\end{equation}

For the free half–line Dirichlet problem:
\begin{equation}
\psi_{\text{free}}
\sim
\sin(kx).
\end{equation}

Matching phases yields
\begin{equation}
\delta
=
-\frac{\pi}{2}\left(\nu-\frac{1}{2}\right).
\end{equation}

Thus the inverse–square potential produces
a purely repulsive phase shift.

\end{document}
